% Generated draft: requires mathematical and source review.
\documentclass[11pt]{article}
\usepackage[margin=25mm]{geometry}
\usepackage{fontspec}
\setmainfont{Latin Modern Roman}
\usepackage{amsmath,amssymb,mathtools,unicode-math}
\setmathfont{Latin Modern Math}
\usepackage{enumitem,needspace,xurl,hyperref}
\setlength{\parindent}{0pt}
\setlength{\parskip}{5pt}
\setlength{\emergencystretch}{3em}
\newcommand{\N}{\mathbb{N}}
\newcommand{\Z}{\mathbb{Z}}
\newcommand{\Q}{\mathbb{Q}}
\newcommand{\R}{\mathbb{R}}
\newcommand{\C}{\mathbb{C}}
\newcommand{\F}{\mathbb{F}}
\newcommand{\A}{\mathbb{A}}
\newcommand{\PP}{\mathbb{P}}
\newcommand{\E}{\mathbb{E}}
\newcommand{\eps}{\varepsilon}
\newcommand{\dd}{\,\mathrm{d}}
\begin{document}
\section*{UnsolvedMath ID: 11000297}
\textbf{Source problem code:} catalogue entry
\subsection{Definitions and mathematical statement}
The question asks for simple Steinberg-type defining relations for the mapping class group. A complete proposal must specify generators and an explicit relation family whose normal closure is the full kernel of the map from the free group on those generators to the intended mapping class group. The source does not define ``Steinberg-type,'' ``simple,'' the generators or surface range, so an arbitrary known finite presentation is not automatically an answer, nor is a particular commutator formula invented.

The following comparison concerns \(\operatorname{Aut}(F_n)\), not the mapping class group. The source describes its action on the Auter-space subcomplex spanned by marked graphs of degree at most two, says this complex is simply connected [48], and applies Brown's method [22] to a presentation. In a usual based-graph degree convention, degree is \(\sum_{v\ne p}(\operatorname{valence}(v)-2)\), rather than a bound on every vertex valence; the precise source convention/rank range is not supplied.

Vertex groups are finite stabilizers of marked graphs. Edge groups stabilize pairs consisting of a marked graph and a forest. Computing the quotient complex gives generation by stabilizers of seven particular marked graphs. All but two of the relations are reported to come from edge-stabilizer inclusions into the stabilizer of the graph or of the graph obtained by collapsing the forest. The seven graphs, quotient complex and two exceptional relations are absent, so no presentation is reconstructed numerically from this description.

The source emphasizes that the group is almost, but not exactly, a pushout of these finite groups: the associated complex of groups is not simple in Haefliger's sense [19, II.12]. This technical simplicity issue belongs to the supplied comparison and need not mean the same thing as the informal request for simple Steinberg-type relations. The two exceptions must not be discarded to turn the report into a pure amalgam/pushout theorem. How this comparison should produce the requested mapping-class relations is itself unstated.
\subsection{Short English statement}
Find a genuinely Steinberg-style presentation of the intended mapping class group. The source compares a seven-stabilizer presentation of \(\operatorname{Aut}(F_n)\), expressly retaining two relations beyond the natural inclusion data.
\subsection{Sources}
\begin{itemize}
\item Local UnsolvedMath registry, permanent ID \(11000297\), exact source code \texttt{catalogue entry}.
\item Farb, \emph{Problems on Mapping Class Groups and Related Topics}, Question 6.1, PDF page 338: \url{https://www.math.uchicago.edu/~farb/papers/mcgbook.pdf}. Registry access 2026-07-29; pdf-font-structure extraction.
\item Internal references [48] (degree-subcomplex simple connectivity), Brown [22] (presentation method), and Haefliger terminology via [19], II.12. Full bibliographic entries are absent.
\end{itemize}
\end{document}
